\label{ch:zustand_prozess}
% Source document: 247

\epigraph{The mere storage of information is energy-neutral.
Only processing has an energy dimension.}{Doc.\,247}

\section{The Question}

In information physics --- represented by Wheeler (``It from Bit''),
M.\,M.~Vopson \textbf{[V1--V3]} and others --- information is often treated as an ontological
primitive, from which energy and matter follow. Landauer's principle is
cited as evidence: Since erasing a bit costs energy,
information is physically real.

This conclusion confuses two structurally different concepts.

\begin{keyresult}
\textbf{Information as state} --- a configuration, a stored bit,
a winding number --- exists without energy expenditure. It is a topological
or geometric property.\\[4pt]
\textbf{Information as process} --- reading, writing, erasing, transforming
--- costs energy. Landauer's principle applies \emph{only} to erasure, thus
to a process, not to the state itself.
\end{keyresult}

\section{Empirical Evidence: The Permanent Memory}

Anyone who flattens the difference must accept a direct empirical consequence:
Every stored bit on a switched-off device would
continuously consume energy --- which would make permanent memory fundamentally
impossible. The existence of flash memory, magnetic hard drives
and ROM chips refutes the equation of information and energy.

A permanent memory holds millions of bits over years at nearly zero
energy consumption. It is a counter-experiment to any theory that equates information
and energy.

\section{The FFGFT Reading}

In FFGFT, states are geometric configurations on $T^4$:
winding numbers $(n_\theta, n_\varphi) \in \mathbb{Z}^2$. Energy is a
derived quantity that follows from \emph{processes} --- from the transition between
configurations --- not from the mere existence of a configuration.

The rest mass of a particle in FFGFT is:
\[
  m = \frac{\hbar c}{L} \cdot n \qquad (n \in \mathbb{Z},\; L \text{ scale})
\]
The winding number $n$ is a state --- it costs no energy as long as
the topology remains intact. Only the transition $n \to n'$ --- a process ---
requires or generates energy.

\section{M.\,M.~Vopson and the FFGFT Restriction}

M.\,M.~Vopson postulates in \textbf{[V1]} a mass-information equivalence: Every bit
carries a universal rest mass. That is attractive --- and according to FFGFT false.

FFGFT restriction: Vopson's bit mass would be system-invariant. FFGFT shows
that the energy of a bit $\Ebit = \hbar c / L$ depends on the scale $L$
(Doc.\,257). Bits on different scales have different energies.
There is no universal bit mass --- neither according to Landauer (Doc.\,A271--A272)
nor according to FFGFT.

\begin{ffgftbox}
\textbf{Wheeler:} ``It from Bit'' --- geometry from information.\\
\textbf{M.\,M.~Vopson \textbf{[V1,V2]}:} Bit mass universal --- information has rest mass.\\
\textbf{FFGFT:} ``Bit from It'' --- information as projection from geometry;
bit mass scale-sensitive; storing is free, erasing costs.
\end{ffgftbox}

\section{What Landauer Really Shows}

Landauer's principle is, in this perspective, not evidence for the
physical reality of information as an ontological primitive.
It is evidence for the physical reality of the \emph{carrier}:
When a carrier is irreversibly reconfigured, costs arise.
The information on the carrier plays no role in this ---
thermodynamics counts regions, not meanings (Doc.\,A271).